% GENERATED FILE. Edit canonical lessons, then run npm run generate:books.
% canonical-source-hash: fnv1a64:09ed8f25c1cac8c4
% canonical-lessons: HI-C414-gusalkhana, HI-C414-saucalay, HI-C414-gizar, HI-C414-aramdeh, HI-C414-jagana

\chapter{Bathrooms and Comfort}
\label{ch:hi-bathrooms-and-comfort}
\hlchaptermodality{414}

\begin{chapteropening}
\textbf{By the end of this chapter:} \emph{I can ask for the bathroom and hot water, say a bed is comfortable and ask to be woken.}

Complete the last lesson of chapter 414: I can ask for the bathroom and hot water, say a bed is comfortable and ask to be woken.
\end{chapteropening}

\section[\texorpdfstring{\hi{गुसलख़ाना} (gusalkhānā)}{गुसलख़ाना (gusalkhānā)}]{\hi{गुसलख़ाना} (gusalkhānā) — a bathroom}

\label{lesson:HI-C414-gusalkhana}



\begin{quote}
Before the new one: say the Hindi for to stop over, then the Hindi for to spread out.
\end{quote}

\subsection*{You'll want to know: \hi{गुसलख़ाना}}
\textbf{\hi{गुसलख़ाना}} — \emph{gusalkhānā} — \textquotedblleft{}a bathroom\textquotedblright{}.

\begin{quote}
\hi{हमारे} \hi{कमरे} \hi{का} \hi{गुसलख़ाना} \hi{छोटा} \hi{है।}
\end{quote}

\begin{quote}
\emph{hamāre kamre kā gusalkhānā choṭā hai.}
\end{quote}

\begin{quote}
The bathroom in our room is small.
\end{quote}

Masculine.

\begin{cousinweb}
\textbf{\hi{गुसलख़ाना}} (\emph{gusalkhānā}) is usually traced to Persian \emph{ġusl-khāna}, a room (\emph{khāna}) for washing (\emph{ġusl}); Hindi speech smooths ġusl to gusal.
\end{cousinweb}

\begin{grammarlens}[title={What you've built}]
One more word for getting around.
\end{grammarlens}

\subsection*{Your turn}
\begin{itemize}
\raggedright
  \item \emph{Say it:} \emph{gusalkhānā}
  \item \emph{Say it:} \emph{gusalkhānā}, once more
  \item \emph{Say it:} \emph{hamāre kamre kā gusalkhānā choṭā hai.} — The bathroom in our room is small.
  \item \emph{Recall:} say the Hindi for to stop over, then the Hindi for to spread out, then say \emph{gusalkhānā} again
\end{itemize}

\begin{tcolorbox}[breakable,colback=teal!4,colframe=teal!35!black,title={Before you move on}]
What does \hi{गुसलख़ाना} mean? (\textquotedblleft{}A bathroom\textquotedblright{}.) Say it once more.
\end{tcolorbox}
\section[\texorpdfstring{\hi{शौचालय} (śaucālay)}{शौचालय (śaucālay)}]{\hi{शौचालय} (śaucālay) — a toilet (formal)}

\label{lesson:HI-C414-saucalay}



\begin{quote}
Before the new one: say the Hindi for to spread out, then the Hindi for a bathroom.
\end{quote}

\subsection*{You'll want to know: \hi{शौचालय}}
\textbf{\hi{शौचालय}} — \emph{śaucālay} — \textquotedblleft{}a toilet (formal)\textquotedblright{}.

\begin{quote}
\hi{स्टेशन} \hi{पर} \hi{शौचालय} \hi{कहाँ} \hi{है}?
\end{quote}

\begin{quote}
\emph{sṭeśan par śaucālay kahā̃ hai?}
\end{quote}

\begin{quote}
Where is the toilet at the station?
\end{quote}

Masculine.

\begin{cousinweb}
\textbf{\hi{शौचालय}} (\emph{śaucālay}) is a learned compound of Sanskrit \emph{śauca}, \textquotedblleft{}cleanness, purification\textquotedblright{}, and \emph{ālaya}, \textquotedblleft{}a house, a dwelling\textquotedblright{} (Monier-Williams); it is the word on public signs.
\end{cousinweb}

\begin{grammarlens}[title={What you've built}]
One more word for getting around.
\end{grammarlens}

\subsection*{Your turn}
\begin{itemize}
\raggedright
  \item \emph{Say it:} \emph{śaucālay}
  \item \emph{Say it:} \emph{śaucālay}, once more
  \item \emph{Say it:} \emph{sṭeśan par śaucālay kahā̃ hai?} — Where is the toilet at the station?
  \item \emph{Recall:} say the Hindi for to spread out, then the Hindi for a bathroom, then say \emph{śaucālay} again
\end{itemize}

\begin{tcolorbox}[breakable,colback=teal!4,colframe=teal!35!black,title={Before you move on}]
What does \hi{शौचालय} mean? (\textquotedblleft{}A toilet (formal)\textquotedblright{}.) Say it once more.
\end{tcolorbox}
\section[\texorpdfstring{\hi{गीज़र} (gīzar)}{गीज़र (gīzar)}]{\hi{गीज़र} (gīzar) — a water heater, a geyser}

\label{lesson:HI-C414-gizar}



\begin{quote}
Before the new one: say the Hindi for a bathroom, then the Hindi for a toilet.
\end{quote}

\subsection*{You'll want to know: \hi{गीज़र}}
\textbf{\hi{गीज़र}} — \emph{gīzar} — \textquotedblleft{}a water heater, a geyser\textquotedblright{}.

\begin{quote}
\hi{गीज़र} \hi{ठीक} \hi{से} \hi{नहीं} \hi{चल} \hi{रहा} \hi{था।}
\end{quote}

\begin{quote}
\emph{gīzar ṭhīk se nahīṃ cal rahā thā.}
\end{quote}

\begin{quote}
The water heater was not working properly.
\end{quote}

Masculine.

\begin{cousinweb}
\textbf{\hi{गीज़र}} (\emph{gīzar}) is borrowed from English \emph{geyser}, which in India names the electric water heater in a bathroom.
\end{cousinweb}

\begin{grammarlens}[title={What you've built}]
One more word for getting around.
\end{grammarlens}

\subsection*{Your turn}
\begin{itemize}
\raggedright
  \item \emph{Say it:} \emph{gīzar}
  \item \emph{Say it:} \emph{gīzar}, once more
  \item \emph{Say it:} \emph{gīzar ṭhīk se nahīṃ cal rahā thā.} — The water heater was not working properly.
  \item \emph{Recall:} say the Hindi for a bathroom, then the Hindi for a toilet, then say \emph{gīzar} again
\end{itemize}

\begin{tcolorbox}[breakable,colback=teal!4,colframe=teal!35!black,title={Before you move on}]
What does \hi{गीज़र} mean? (\textquotedblleft{}A water heater, a geyser\textquotedblright{}.) Say it once more.
\end{tcolorbox}
\section[\texorpdfstring{\hi{आरामदेह} (ārāmdeh)}{आरामदेह (ārāmdeh)}]{\hi{आरामदेह} (ārāmdeh) — comfortable}

\label{lesson:HI-C414-aramdeh}



\begin{quote}
Before the new one: say the Hindi for a toilet, then the Hindi for a water heater.
\end{quote}

\subsection*{You'll want to know: \hi{आरामदेह}}
\textbf{\hi{आरामदेह}} — \emph{ārāmdeh} — \textquotedblleft{}comfortable\textquotedblright{}.

\begin{quote}
\hi{होटल} \hi{का} \hi{बिस्तर} \hi{बहुत} \hi{आरामदेह} \hi{था।}
\end{quote}

\begin{quote}
\emph{hoṭal kā bistar bahut ārāmdeh thā.}
\end{quote}

\begin{quote}
The hotel bed was very comfortable.
\end{quote}

\begin{cousinweb}
\textbf{\hi{आरामदेह}} (\emph{ārāmdeh}) is usually traced to Persian \emph{ārāmdih}, \textquotedblleft{}giving rest\textquotedblright{}, from \emph{ārām}, \textquotedblleft{}rest\textquotedblright{}, the word in \hi{आराम} \hi{से}. It does not change for gender.
\end{cousinweb}

\begin{grammarlens}[title={What you've built}]
One more word for getting around.
\end{grammarlens}

\subsection*{Your turn}
\begin{itemize}
\raggedright
  \item \emph{Say it:} \emph{ārāmdeh}
  \item \emph{Say it:} \emph{ārāmdeh}, once more
  \item \emph{Say it:} \emph{hoṭal kā bistar bahut ārāmdeh thā.} — The hotel bed was very comfortable.
  \item \emph{Recall:} say the Hindi for a toilet, then the Hindi for a water heater, then say \emph{ārāmdeh} again
\end{itemize}

\begin{tcolorbox}[breakable,colback=teal!4,colframe=teal!35!black,title={Before you move on}]
What does \hi{आरामदेह} mean? (\textquotedblleft{}Comfortable\textquotedblright{}.) Say it once more.
\end{tcolorbox}
\section[\texorpdfstring{\hi{जगाना} (jagānā)}{जगाना (jagānā)}]{\hi{जगाना} (jagānā) — to wake (someone)}

\label{lesson:HI-C414-jagana}



\begin{quote}
Before the new one: say the Hindi for a water heater, then the Hindi for comfortable.
\end{quote}

\subsection*{You'll want to know: \hi{जगाना}}
\textbf{\hi{जगाना}} — \emph{jagānā} — \textquotedblleft{}to wake (someone)\textquotedblright{}.

\begin{quote}
\hi{मुझे} \hi{सुबह} \hi{पाँच} \hi{बजे} \hi{जगाइए।}
\end{quote}

\begin{quote}
\emph{mujhe subah pā̃c baje jagāiye.}
\end{quote}

\begin{quote}
Please wake me at five in the morning.
\end{quote}

Past, with \hi{ने} (he or she …): \hi{उसने} \hi{मुझे} \hi{जल्दी} \hi{जगाया}, \emph{usne mujhe jaldī jagāyā}.

\begin{cousinweb}
\textbf{\hi{जगाना}} (\emph{jagānā}) is the causative partner of \hi{जागना}, from Sanskrit \emph{jāgrati}, \textquotedblleft{}are awake\textquotedblright{} (CDIAL 5175), which lists Hindi \emph{jagānā}: you \hi{जागना} yourself, and you \hi{जगाना} someone else.
\end{cousinweb}

\begin{grammarlens}[title={What you've built}]
One more word for getting around.
\end{grammarlens}

\subsection*{Your turn}
\begin{itemize}
\raggedright
  \item \emph{Say it:} \emph{jagānā}
  \item \emph{Say it:} \emph{jagānā}, once more
  \item \emph{Say it:} \emph{mujhe subah pā̃c baje jagāiye.} — Please wake me at five in the morning.
  \item \emph{Recall:} say the Hindi for a water heater, then the Hindi for comfortable, then say \emph{jagānā} again
\end{itemize}

\begin{tcolorbox}[breakable,colback=teal!4,colframe=teal!35!black,title={Before you move on}]
What does \hi{जगाना} mean? (\textquotedblleft{}To wake (someone)\textquotedblright{}.) Say it once more.
\end{tcolorbox}
